\begin{frame}
	\frametitle{Duas fronteiras}
	\small
	\begin{columns}
		\column{.48\linewidth}
		\begin{beamercolorbox}[sep=5pt,rounded=true,shadow=true]{titleBitnet}
			\textbf{Fronteira 1 --- BitNets}
		\end{beamercolorbox}
		\vspace{1pt}
		\par Eixo da \textbf{representação}: quantos bits descrevem cada peso/ativação.
		\begin{itemize}
			\item FP32 $\rightarrow$ FP16 $\rightarrow$ INT8 $\rightarrow$ ternário $\{-1,0,1\}$ $\rightarrow$ binário $\{-1,1\}$.
			\item A rede continua densa e síncrona: \textbf{todo} neurônio ainda computa em \textbf{toda} passada.
			\item O que muda é o \textbf{custo} de cada computação individual.
		\end{itemize}
		\column{.48\linewidth}
		\begin{beamercolorbox}[sep=5pt,rounded=true,shadow=true]{titleSnn}
			\textbf{Fronteira 2 --- Redes de pulso}
		\end{beamercolorbox}
		\vspace{1pt}
		\par Eixo do \textbf{tempo/atividade}: quando e quanto se computa.
		\begin{itemize}
			\item Ativação contínua e síncrona $\rightarrow$ pulso esparso, assíncrono, orientado a evento.
			\item Os pesos podem continuar em precisão plena.
			\item O que muda é \textbf{se} um neurônio computa, a cada instante.
		\end{itemize}
	\end{columns}
	\begin{block}{O alvo é o mesmo}
		Ambas atacam o custo (energia, memória, latência) da multiplicação densa em ponto flutuante --- mas por caminhos diferentes.
	\end{block}
\end{frame}
